% Imported from LaTeX_and_pdfs/AGAT_dump.tex; see provenance.json.
\part{Waves}


\section{Forced Oscillation and Resonance, and Damped Oscillations}

See Georgi (1992) \cite{agat:Geor1992}.

\subsection{Damped Oscillators}

See Sec. 2.1 "Damped Oscillations", Georgi (1992) \cite{agat:Geor1992}.

For small velocities, frictional force takes the form of

\begin{equation}
	-m\Gamma v
\end{equation}
See (2.1.1) of Georgi (1992) \cite{agat:Geor1992}. $1/\Gamma$ has the dimensions of time.

Now
\[
\begin{gathered}
	\sum F = m\ddot{x} = -m \Gamma v - Kx
\end{gathered}
\]
or
\begin{equation}
	\ddot{x} + \Gamma \dot{x} + \frac{K}{m} x = \ddot{x} + \Gamma \dot{x} + \omega_0^2 x = 0
\end{equation}
where $\omega_0 = \sqrt{ \frac{K}{m} }$. See (2.1.2) of Georgi (1992) \cite{agat:Geor1992}.

Since $\ddot{x} + \Gamma \dot{x} + \omega_0^2 x = 0$ is 1. time translation invariant, i.e. if $x(t+a)$ is a solution, then $x(t)$ is a solution, and 2. linearity, then find irreducible solutions of form $z(t) = e^{\alpha t}$.

\[
\Longrightarrow \alpha^2 + \Gamma \alpha + \omega_0^2 =0 \text{ or } \alpha = \frac{ -\Gamma \pm \sqrt{ \Gamma^2 - 4\omega_0^2}}{2}
\]

\subsection{Overdamped oscillations} See Sec. 2.1.1 "Overdamped Oscillations" of Georgi (1992) \cite{agat:Geor1992}.

If $\frac{\Gamma}{2} > \omega_0$, let
\begin{equation}
	\Gamma_{\pm} = \frac{\Gamma}{2} \pm \sqrt{ \frac{\Gamma^2}{4} - \omega_0^2}
\end{equation}
See (2.1.8) of Georgi (1992) \cite{agat:Geor1992}

\begin{equation}
	x(t) = A_+ e^{-\Gamma_+ t} + A_- e^{-\Gamma_- t}
\end{equation}
See (2.1.7) of Georgi (1992) \cite{agat:Geor1992}.

\subsection{Underdamped Oscillations} See Sec. 2.1.2 "Underdamped Oscillations" of Georgi (1992) \cite{agat:Geor1992}.

If $\frac{\Gamma}{2} < \omega_0$, so
\[
\alpha = \frac{-\Gamma }{2} \pm i \sqrt{ \omega_0^2 - \frac{\Gamma^2}{2} } 
\]

so 
\begin{equation}
	\omega^2 = \omega_0^2 - \frac{ \Gamma^2}{4}
\end{equation}

\begin{equation}
	x(t) = Ae^{\frac{-\Gamma}{2} t} \cos{ (\omega t- \theta)}
\end{equation}
See (2.1.10) of Georgi (1992) \cite{agat:Geor1992}.

\[
x(t) = A e^{\frac{-\Gamma}{2} t} c(\omega t - \theta) = Ae^{\frac{-\Gamma}{2}t } (c(\omega t) c(\theta) + s(\omega t) s(\theta))
\]
\begin{equation}\label{agat:Eq:UnderdampedOscillationsSolution}
	x(t) = e^{-\frac{\Gamma}{2}t} (c \cos{(\omega t)} + d\sin{(\omega t)})
\end{equation}
See (2.1.11) of Georgi (1992) \cite{agat:Geor1992}, where $c= A\cos{\theta}$, $d=A\sin{\theta}$.

\subsubsection{Critically Damped Oscillations} See Sec. 2.1.3 "Critically Damped Oscillations" of Georgi (1992) \cite{agat:Geor1992}.

If $\frac{\Gamma}{2} = \omega_0$, we only get $e^{\frac{-\Gamma}{2} t}$. We know that there will be 2 solutions to the 2nd-order differential equation.

One way to find the other solution is to approach this situation from underdamped case as a limit, i.e.
\[
\omega \to 0 \text{ from } \omega > 0
\]

\[
\begin{gathered}
	\begin{aligned}
		& \lim_{\omega \to 0} e^{ \frac{-\Gamma}{2}(t) } c\cos{(\omega t)} = e^{ \frac{-\Gamma}{2} t} \\
		& \lim_{\omega\to 0} e^{\frac{-\Gamma}{2} t} d \sin{\omega t} = 0
	\end{aligned}
\end{gathered}
\]

Divide the second term by $\omega$, because $\omega$ doesn't depend on $\omega$, $\frac{1}{\omega} e^{\frac{-\Gamma}{2} t} d\sin{\omega t}$ is still a solution.

\[
\begin{gathered}
	\lim_{\omega \to 0} e^{ \frac{-\Gamma}{2} t} \frac{ \sin{ \omega t} }{\omega} = te^{\frac{-\Gamma}{2} t} (\text{by L'Hopital's rule})
\end{gathered}
\]
So
\[
x(t) = ce^{ - \Gamma t/2} + dte^{-\Gamma t/2}
\]

\subsection{Forced oscillations} See Sec. 2.2 "Forced Oscillations" of Georgi (1992) \cite{agat:Geor1992}.
Equation of motion is
\begin{equation}\label{agat:Eq:ForcedOscillationEquationOfMotion}
	\ddot{x}(t) + \Gamma \dot{x} + \omega_0^2 x = F(t) /m 
\end{equation}
where
\begin{equation}
	F(t) = F_0 \cos{(\omega_dt)}
\end{equation}
See (2.1.5), (2.1.4) of Georgi (1992) \cite{agat:Geor1992}, respectively. \\

$\frac{\omega_d}{2\pi} = $ driving frequency, driving frequency can be tuned completely independently of the other parameters of the system.

If we let 
\[
\ddot{z}(t) + \Gamma \dot{z}(t) + \omega_0^2 z(t) = \mathcal{F}(t) /m
\]
where $\mathcal{F}(t) = F_0 e^{-i \omega_0 t}$.

For in particular numerical computation applications, we can transform Eq. \ref{agat:Eq:ForcedOscillationEquationOfMotion} into a system of first order equations: \\
Let 
\[
\begin{aligned}
	& x_1(t) = x(t) \\
	& x_2(t) = \dot{x}(t)
\end{aligned}
\]
so that
\[
\begin{aligned}
	& \dot{x}_1 = x_2 \\
	& \dot{x}_2 = -\Gamma x_2 - \omega_0^2 x_1 + f(t)
\end{aligned}
\]

Then if we look for steady state solution of the form
\[
z(t) = A e^{-i \omega_d t}
\]

One can think like the steady state solution is such that the system is driven with driving frequency $\omega_d$.

\begin{equation}
	(-\omega_d^2 + (-i \omega_d)\Gamma + \omega_0^2)A = \frac{F_0}{m} 
\end{equation}
See (2.20) of Georgi (1992) \cite{agat:Geor1992}

or
\begin{equation}\label{agat:Eq:ForcedOscillationSteadyStateSolutionAmplitude}
	A = \frac{ F_0 / m }{ \omega_0^2 - i\omega_d \Gamma - \omega_d^2}
\end{equation}
See (2.21) of Georgi (1992) \cite{agat:Geor1992}

The amplitude $A$ of the displacement is proportional to amplitude of driving force; this is as expected from linearity.

\[
\begin{gathered}
	\begin{aligned}
		& A = \frac{ \frac{F_0}{m} (\omega_0^2 + i \omega_d \Gamma - \omega_d^2 )}{ (\omega_0^2 - \omega_d^2)^2 + (\omega_d \Gamma)^2 } \text{ with } & \\
		& \Re{A} = \frac{ (\omega_0^2 - \omega_d^2) F_0/m }{(\omega_0^2 - \omega_d^2)^2 + (\omega_d \Gamma)^2} & \Im{A} \equiv B = \frac{ \Gamma \omega_d F_0 /m }{ (\omega_0^2 - \omega_d^2)^2 + (\Gamma \omega_d )^2 }
	\end{aligned} \\
	\mathcal{A} e^{- i\omega_d t} = \mathcal{A} ( c(\omega_dt) - i s(\omega_dt)) \xrightarrow{\Re} \Re{A } c(\omega_d t) + \Im{A} s(\omega_d t)
\end{gathered}
\]
where I use $\mathcal{A} \equiv A$ notation interchangeably, relying on context on whether it's real or complex, and $\cos{\omega_d t} \equiv c(\omega_d t)$, $\sin{\omega_d t} \equiv s(\omega_d t)$ as shorthand notation for sin and cos. \\

For 
\[
x(t) = \Re{z(t)} = (\Re{A}) c(\omega_d t) + Bs(\omega_d t)
\]

Term proportional to $B$ is $90^{\circ}$ out of phase. \\
Real part of $\mathcal{A}$, $A = \Re{A}$ is the \emph{elastic amplitude.} \\
Imaginary part of $\mathcal{A}$, $B = \Im{A}$ is called the \emph{absorptive amplitude}. \\

In the case where 
\[
F(t) = F_0 \sin{(\omega_dt)}
\]
we'll see that the fundamental physical characteristics don't change, but the initial conditions for the solution do change, which is important for actual modeling (but gets looked over a bit in physics that concern about high-level concepts).

If $F(t) = F_0 \sin{(\omega_d t)}$, while the following doesn't change:
\[
\begin{gathered}
	\begin{aligned}
		& z(t) = \mathcal{A} e^{-i\omega_d t} \\
		& \dot{z}(t) = (-i \omega_d) z \\
		& \ddot{z}(t) = (-\omega_d^2)z
	\end{aligned}
\end{gathered}
\]
but now consider 
\[
\begin{gathered}
	\mathcal{F}(t) = + i F_0 e^{-i\omega_d t} = +i F_0 (c(\omega_d t)  -i s(\omega_d t) ) \\
	\Re{(\mathcal{F}(t) )} = F_0 s(\omega_dt)
\end{gathered}
\]
so then
\[
\begin{gathered}
	(-\omega_d^2 + (-i \omega_d)\Gamma + \omega_0^2)\mathcal{A} = \frac{iF_0}{m}  \text{ or } 	\\
	\mathcal{A} = \frac{ iF_0 / m }{ \omega_0^2 - \omega_d^2- i\omega_d \Gamma} = \frac{ \frac{iF_0}{m} (\omega_0^2 - \omega_d^2 + i \omega_d \Gamma  )}{ (\omega_0^2 - \omega_d^2)^2 + (\omega_d \Gamma)^2 }
\end{gathered}
\]

So then in this case, 
\begin{equation}\label{agat:Eq:AmplitudeValuesForcedOscillationsForceProportionalToSine}
	\begin{aligned}
		& \Re{A} = \frac{ -\Gamma \omega_d F_0 /m }{ (\omega_0^2 - \omega_d^2)^2 + (\Gamma \omega_d )^2 } \equiv A \\
		& \Im{A} = \frac{ (\omega_0^2 - \omega_d^2) F_0/m }{(\omega_0^2 - \omega_d^2)^2 + (\omega_d \Gamma)^2} \equiv B 
	\end{aligned}
\end{equation}

and so the general form of the steady-state solution doesn't change, but the values for the amplitudes get switched:

\[
\begin{gathered}
	x(t) = \Re{ z(t)} = Ac(\omega_d t) + Bs(\omega_dt) \\
	\dot{x}(t) = \omega_d \left[ -A s(\omega_dt) + Bc(\omega_dt) \right]
\end{gathered}
\]

\subsection{Resonances} See Sec. 2.3 "Resonances" of Georgi (1992) \cite{agat:Geor1992}.

\[
\omega_d = \omega_0
\]
so $\Re{A} = 0$, $B = \frac{ F_0 / m}{ \Gamma \omega_d}$.

\subsection{Work, Work done by forced oscillation on the system}

Work is a nonlinear function of force, so we must use real force and real displacement.

\[
P(t) = F(t) \frac{\partial }{\partial t} X(t) = F_0 c(\omega_d t) (\omega_d) (-\Re{A} s(\omega_d t) + B c(\omega_d t))
\]

The first term is proportional to $s(2\omega_d t)$ and averages to 0 over a half period.

\[
\int_{t_0}^{t_0 + \pi / \omega_d} s(2\omega_d t) dt = \left. \frac{ -c(2\omega_d t) }{2 \omega_d } \right|_{t_0}^{t_0 + \pi /\omega_d} = 0
\]

This is why $\Re{A}$ is called elastic amplitude. If $A$ dominates, then energy fed into the system at 1 time is returned at a later time, as in elastic collision in mechanics.
\[
\begin{gathered}
	P_{\text{average}} \text{of 2nd. term} = P_{\text{average}} \\
	\text{ Consider } \\
	F_0 \omega_d B \int_{t_0}^{t_0 + \pi/ \omega_d} c^2 (\omega_d t) dt = F_0 \omega_d B \int_{t_0}^{t_0 + \pi/ \omega_d} \frac{ c(2\omega_d t) + 1}{2} dt = \\
	F_0 \omega_d B \left. \left( \frac{ s(2\omega_d t) }{ 4 \omega_d } + \frac{t}{2} \right) \right|_{t_0}^{t_0 + \frac{\pi}{\omega_d} } = F_0 \omega_d B \frac{\pi}{2\omega_d} = \frac{F_0 B\pi}{2} \\
	P_{\text{average}} = \frac{ \int_{t_0}^{t_0 + \pi / \omega_d} P(t) }{ \Delta t}  \text{ where } \Delta t = \pi / \omega_d \\
	\Longrightarrow P_{\text{average}} = \frac{ F_0 B\pi /2}{ \pi / \omega_d}  = \boxed{ \frac{F_0 \omega_d B}{2} }
\end{gathered}
\]

Recall that
\[
B = \frac{ \Gamma \omega_d F_0 /m }{ (\omega_0^2 - \omega_d^2)^2 + (\Gamma \omega_d)^2 }
\]

\[
\begin{gathered}
	P_{\text{average}} = \frac{F_0 \omega_d B}{2} = \frac{ (F_0 \omega_d)^2 \Gamma /2m }{ (\omega_0^2 - \omega_d^2)^2 + (\Gamma \omega_d)^2 } \xrightarrow{ \omega_d\to \omega_0} \frac{F_0^2}{ 2m \Gamma}
\end{gathered}
\]

The maximum average power is inversely proportional to $\Gamma$. \\

This analysis doesn't change when $F(t) = F_0 \sin{(\omega_d t)}$.

\[
\begin{gathered}
	P(t) = F(t) \dot{x}(t) = F_0 s(\omega_d t) \omega_d [ -A s(\omega_d t) + Bc(\omega_d t) ] \\
	\begin{aligned}
		& \int_{t_0}^{t_0 + \frac{\pi}{\omega_d} } dt c(\omega_dt) s(\omega_d t) = \left. \frac{1}{2 \omega_d} s^2(\omega_d t) \right|_{t_0}^{t_0 + \frac{\pi}{\omega_d} } = 0 \\
		& \int_{t_0}^{t_0 + \frac{\pi}{\omega_d} } dt s^2(\omega_dt) = \int_{t_0}^{t_0 + \frac{\pi}{\omega_d} } dt \frac{1 - c(2\omega_dt)}{2} = \frac{\pi}{2 \omega_d} - \left. \frac{s(2\omega_dt)}{ 4\omega_d} \right|_{t_0}^{t_0 + \frac{\pi}{\omega_d} } = \frac{\pi}{2\omega_d}
	\end{aligned} \\
	P_{\text{average}} = \frac{ \int_{t_0}^{t_0 + \frac{\pi}{\omega_d} } dt P(t) }{ \pi /\omega_d} = \frac{-F_0 A \pi}{2 / \pi/\omega_d} = \frac{-F_0 A \omega_d}{2}
\end{gathered}
\]
So $B$ is now the elastic amplitude and $A$ is the absorptive amplitude, but the same values are obtained, since $A$ and $B$ are effectively switched for $F(t) = F_0 \sin{\omega_d t}$.

To find the width (for fixed height) of the resonance curve (curve of $P_{\text{average}}$, depending on $\omega_d$), \\
find the value $\omega_d = \omega_{1/2}$ s.t. 
\[
P_{\text{average}} = \frac{P_{\text{average,max}}}{2} = \frac{F_0^2}{4m \Gamma}
\]
Then

\[
\begin{gathered}
	\frac{P_{\text{average,max}}}{2} = \frac{F_0^2}{4m \Gamma} = \frac{ (F_0^2 \Gamma /2m ) \omega_{1/2}^2 }{ (\omega_{1/2}^2 - \omega_0^2)^2 + (\Gamma \omega_{1/2})^2 } \text{ or } \omega_{1/2}^4 -2\omega_{1/2}^2 \omega_0^2 + \omega_0^2 + \Gamma^2 \omega_{1/2}^2 = 2\Gamma^2 \omega_{1/2}^2 \text{ or } \\
	\begin{aligned}
		& \omega_{1/2}^4 - (2\omega_0^2 + \Gamma^2) \omega_{1/2}^2 + \omega_0^4 = 0 \\
		& \omega_{1/2}^2 = \frac{ (2 \omega_0^2 + \Gamma^2) \pm \sqrt{ ( \Gamma^2 + 2\omega_0^2)^2 - 4\omega_0^4 } }{2} = \omega_0^2 + \frac{ \Gamma^2}{2} \pm \sqrt{ \frac{\Gamma^4 }{4}  + \Gamma^2 \omega_0^2 } = \omega_0^2 + \frac{ \Gamma^2}{2} \pm \Gamma \sqrt{ \frac{\Gamma^2 }{4}  +  \omega_0^2 }
	\end{aligned}
\end{gathered}
\]
Now (EY: it seems like you would have needed to have known the answer beforehand to have tried this factorization)
\[
\begin{gathered}
	\omega_{1/2}^2 = \left( \sqrt{ \omega_0^2 + \frac{\Gamma^2}{4} } \right)^2 \pm 2 \frac{\Gamma}{2} \sqrt{ \frac{\Gamma^2}{4} + \omega_0^2 } + \left( \frac{\Gamma}{2} \right)^2  = \left( \sqrt{ \omega_0^2 + \frac{ \Gamma^2}{4} } \pm \frac{\Gamma}{2} \right)^2
\end{gathered}
\]
Then
\begin{equation}
	\boxed{ \omega_{1/2} = \sqrt{ \omega_0^2 + \frac{\Gamma^2}{4} } \pm \frac{\Gamma}{2} }
\end{equation}

The "free oscillation" of the system would be the case of underdamped oscillations; recall
\[
x(t) = A \exp{ \left( \frac{-\Gamma t}{2} \right) }\cos{ (\omega t - \theta) }
\]
Notice the exponential decay behavior of the amplitude or i.e. magnitude of $x(t)$ due to the $\exp{ \left( \frac{-\Gamma t}{2} \right)}$ factor. Let's just say the "lifetime" of free oscillation is the time for the amplitude to drop to half its value, i.e.
\[
\begin{gathered}
	\frac{1}{2} = \exp{ \left( \frac{-\Gamma t_{1/2} }{2} \right) } \text{ or } -\ln{2} = \frac{-\Gamma t_{1/2} }{2} , \quad \text{ or } t_{1/2} = \frac{2\ln{2} }{\Gamma}
\end{gathered}
\]

For the time dependence of the free oscillation, i.e. $x(t) = A \exp{ \left( \frac{-\Gamma t}{2} \right) }\cos{ (\omega t - \theta) }$, then we say that the lifetime of the state in free oscillation is of order $1/\Gamma$, i.e.
\[
t_{1/2} = \frac{2\ln{2}}{ \Gamma}
\]

The width of the resonance peak in forced oscillation is 
\[
\Delta \omega_{1/2} = \Gamma
\]
so width of the resonance peak in forced oscillation is inversely proportional to the lifetime of the corresponding normal mode of free oscillation.

\subsection{Apply Fourier Transform to Forced Oscillations}

So for the equation of motion, Eq. \ref{agat:Eq:ForcedOscillationEquationOfMotion},
\[
\ddot{x} + \Gamma \dot{x} + \omega_0^2 x = F(t)
\]

Assume we drive the system at frequency $\omega_d$ and that system's response's frequency is also $\omega_d$.

Then

\[
\begin{gathered}
	(i\omega_d)^2 \widehat{x} + (i\omega_d) \Gamma \widehat{x} + \omega_0^2 \widehat{x} = \widehat{F}(\omega_d) \text{ or } \frac{ \widehat{x}(\omega_d) }{ \widehat{F}(\omega_d) } = \frac{1}{ - (\omega_d^2 - i \omega_d \Gamma - \omega_0^2) }
\end{gathered}
\]

So in general, $\frac{ \widehat{x}(\omega_d) }{ \widehat{F}(\omega_d) }$ has a real part and imaginary part.


\subsection{Apply Laplace Transform to Forced Oscillations}

Recall Eq. \ref{agat:Eq:ForcedOscillationEquationOfMotion},
\[
\ddot{x} + \Gamma \dot{x} + k x = mF(t)
\]
so that if $\frac{k}{m} = \omega_0^2$, then
\[
\begin{gathered}
	\xrightarrow{ \mathcal{L}} s^2 \widehat{x}(s) - s x(0) - x'(0) + \Gamma (s\widehat{x}(s) - x(0)) + \omega_0^2 \widehat{x}(s) = \widehat{F}(s) \text{ or } \\
	(s^2 + \Gamma s + \omega_0^2) \widehat{x}(s) - (\Gamma + s) x(0) - x'(0) = \widehat{F}(s)
\end{gathered}
\]

If $x'(0) = 0$, and we start system at $x(0) =0$, so that we rely upon $F(t)$ to drive system, then we get a ratio of output or "response" to input, the driving, forced oscillation, per mass,
\[
\begin{gathered}
	\frac{ \widehat{x}(s) }{ \widehat{F}(s) } = \frac{1}{ s^2 + \Gamma s + \omega_0^2 } = \frac{1}{ \left( s + \frac{ \Gamma + \sqrt{ \Gamma^2 - 4 \omega_0^2 } }{ 2 } \right) \left( s + \frac{ \Gamma - \sqrt{ \Gamma^2 - 4 \omega_0^2 } }{ 2} \right) }
\end{gathered}
\]
Now look at where the denominator "blows up", the singularities. We had the case that if $\Gamma^2 - 4\omega^2 \geq 0 $ or $\frac{\Gamma}{2} \geq \omega_0$, we had overdamped or critically damped exponential decay. Then, the value of $s$ where a singularity occurs would be completely real. For the underdamped case with free oscillation, the value of $s$ where a singularity occurs would have an imaginary part.

If initial conditions can't be ignored, \\
then if $x'(0) =0$,
\[
\begin{gathered}
	(s^2 + \Gamma s + \omega_0^2) \widehat{x}(s) - x'(0) = \widehat{F}(s) \text{ or } \frac{ \widehat{x}(s) }{ \widehat{F}(s) } = \frac{ 1 + x'(0) / \widehat{F}(s) }{ s^2 + \Gamma s + \omega_0^2 }
\end{gathered}
\]

But if $x'(0) =0$,
\[
\begin{gathered}
	\frac{ \widehat{x}(s) }{ \widehat{F}(s) } = \frac{ 1 + (\Gamma +s )x(0) / \widehat{F}(s) }{ s^2 + \Gamma s + \omega_0^2 }
\end{gathered}
\]

\subsection{General solution to forced oscillation}

For subscript $s$ meaning steady-state driven solution, see Eq. \ref{agat:Eq:ForcedOscillationSteadyStateSolutionAmplitude}
\[
\begin{gathered}
	x_s(t) = Ac(\omega_d t) + Bs(\omega_d t) \\
	\begin{aligned}
		& A = \frac{ (\omega_0^2 - \omega_d^2) F_0 / m }{ (\omega_0^2 - \omega_d^2)^2 + (\omega_d \Gamma)^2 } \\
		& B = \frac{ \Gamma \omega_d F_0 / m }{ (\omega_0^2 - \omega_d^2)^2 + (\Gamma \omega_d)^2 }
	\end{aligned}
\end{gathered}
\] 
so then
\[
\begin{aligned}
	& x_s(0) = A \\
	& \dot{x}_s(0) = \omega_dB
\end{aligned}
\]

Let $x_{\text{free}}$ denote the response, or solution, or the resulting motion for the "free oscillation", the underdamped, exponential decaying solution, for $\omega_0 > \frac{\Gamma}{2}$. See Eq. \ref{agat:Eq:UnderdampedOscillationsSolution},

\[
\begin{gathered}
	\begin{aligned}
		x_{\text{free}}(t) & = \exp{ \left( \frac{-\Gamma t}{2} \right) } (c\cos{(\omega t)} + d\sin{ (\omega t)} ), \quad \, \omega^2 = \omega_0^2 -\frac{\Gamma^2}{2} \\
		x_{\text{free}}(0) &= c \\
		\dot{x}_{\text{free}}(t) & = \frac{-\Gamma}{2} x_{\text{free}}(t) + e^{-\frac{\Gamma}{2} t} \omega( - cs(\omega t) + dc(\omega t) ) \\
		\dot{x}_{\text{free}}(0) & = \frac{-\Gamma}{2} c + \omega d
	\end{aligned}
\end{gathered}
\]

The most general solution $x(t)$ is of the form
\[
\begin{gathered}
	x(t) = x_{\text{free}}(t) + x_s(t)
\end{gathered}
\]
with initial conditions
\[
\begin{aligned}
	& x(0) = c + A \text{ so } c = x(0) - A \\
	& \dot{x}(0) = \frac{-\Gamma}{2} c + \omega d + \omega_d B \text{ so } d = \frac{1}{\omega} \left[ \dot{x}(0) + \frac{\Gamma}{2} (x(0) - A) - \omega_d B \right]
\end{aligned}
\]

And so the full, general solution is 
\begin{equation}
	\begin{gathered}
		x(t) = e^{\frac{-\Gamma t}{2} } (c\cos{(\omega t)} + d \sin{ (\omega t)} ) + A \cos{ (\omega_d t)} + B \sin{(\omega_d t)} \\
		\text{ where } \\
		\begin{aligned}
			& A = \frac{ (\omega_0^2 - \omega_d^2) F_0 / m }{ (\omega_0^2 - \omega_d^2)^2 + (\omega_d \Gamma)^2 } \\
			& B = \frac{ \Gamma \omega_d F_0 / m }{ (\omega_0^2 - \omega_d^2)^2 + (\Gamma \omega_d)^2 }
			& \omega^2 = \omega_0^2 - \frac{\Gamma^2}{4}, \\
			& c = x(0) - A, \\
			& d = \frac{1}{\omega} \left[ \dot{x}(0) + \frac{\Gamma}{2} (x(0) - A) - \omega_d B \right]
		\end{aligned}
	\end{gathered}
\end{equation}

For the case of $F(t) = F_0 \sin{(\omega_d t)}$, i.e. the force is proportional to $\sin$, then the exact same analysis leads to the same result as above, but for $A,B$ values switched to be as Eq. \ref{agat:Eq:AmplitudeValuesForcedOscillationsForceProportionalToSine}.

\[
\begin{aligned}
	& \Re{A} = \frac{ -\Gamma \omega_d F_0 /m }{ (\omega_0^2 - \omega_d^2)^2 + (\Gamma \omega_d )^2 } \equiv A \\
	& \Im{A} = \frac{ (\omega_0^2 - \omega_d^2) F_0/m }{(\omega_0^2 - \omega_d^2)^2 + (\omega_d \Gamma)^2} \equiv B 
\end{aligned}
\]

\subsection{Forced Oscillation and Resonance, and Damped Oscillations}

See Georgi (1992) \cite{agat:Geor1992}.

\subsubsection{Damped Oscillators}

See Sec. 2.1 "Damped Oscillations", Georgi (1992) \cite{agat:Geor1992}.

For small velocities, frictional force takes the form of

\begin{equation}
	-m\Gamma v
\end{equation}
See (2.1.1) of Georgi (1992) \cite{agat:Geor1992}. $1/\Gamma$ has the dimensions of time.

Now
\[
\begin{gathered}
	\sum F = m\ddot{x} = -m \Gamma v - Kx
\end{gathered}
\]
or
\begin{equation}
	\ddot{x} + \Gamma \dot{x} + \frac{K}{m} x = \ddot{x} + \Gamma \dot{x} + \omega_0^2 x = 0
\end{equation}
where $\omega_0 = \sqrt{ \frac{K}{m} }$. See (2.1.2) of Georgi (1992) \cite{agat:Geor1992}.

Since $\ddot{x} + \Gamma \dot{x} + \omega_0^2 x = 0$ is 1. time translation invariant, i.e. if $x(t+a)$ is a solution, then $x(t)$ is a solution, and 2. linearity, then find irreducible solutions of form $z(t) = e^{\alpha t}$.

\[
\Longrightarrow \alpha^2 + \Gamma \alpha + \omega_0^2 =0 \text{ or } \alpha = \frac{ -\Gamma \pm \sqrt{ \Gamma^2 - 4\omega_0^2}}{2}
\]

\subsubsection{Overdamped oscillations} See Sec. 2.1.1 "Overdamped Oscillations" of Georgi (1992) \cite{agat:Geor1992}.

If $\frac{\Gamma}{2} > \omega_0$, let
\begin{equation}
	\Gamma_{\pm} = \frac{\Gamma}{2} \pm \sqrt{ \frac{\Gamma^2}{4} - \omega_0^2}
\end{equation}
See (2.1.8) of Georgi (1992) \cite{agat:Geor1992}

\begin{equation}
	x(t) = A_+ e^{-\Gamma_+ t} + A_- e^{-\Gamma_- t}
\end{equation}
See (2.1.7) of Georgi (1992) \cite{agat:Geor1992}.

\subsubsection{Underdamped Oscillations} See Sec. 2.1.2 "Underdamped Oscillations" of Georgi (1992) \cite{agat:Geor1992}.

If $\frac{\Gamma}{2} < \omega_0$, so
\[
\alpha = \frac{-\Gamma }{2} \pm i \sqrt{ \omega_0^2 - \frac{\Gamma^2}{2} } 
\]

so 
\begin{equation}
	\omega^2 = \omega_0^2 - \frac{ \Gamma^2}{4}
\end{equation}

\begin{equation}
	x(t) = Ae^{\frac{-\Gamma}{2} t} \cos{ (\omega t- \theta)}
\end{equation}
See (2.1.10) of Georgi (1992) \cite{agat:Geor1992}.

\[
x(t) = A e^{\frac{-\Gamma}{2} t} c(\omega t - \theta) = Ae^{\frac{-\Gamma}{2}t } (c(\omega t) c(\theta) + s(\omega t) s(\theta))
\]
\begin{equation}\label{agat:Eq:UnderdampedOscillationsSolution:2}
	x(t) = e^{-\frac{\Gamma}{2}t} (c \cos{(\omega t)} + d\sin{(\omega t)})
\end{equation}
See (2.1.11) of Georgi (1992) \cite{agat:Geor1992}, where $c= A\cos{\theta}$, $d=A\sin{\theta}$.

\subsubsection{Critically Damped Oscillations} See Sec. 2.1.3 "Critically Damped Oscillations" of Georgi (1992) \cite{agat:Geor1992}.

If $\frac{\Gamma}{2} = \omega_0$, we only get $e^{\frac{-\Gamma}{2} t}$. We know that there will be 2 solutions to the 2nd-order differential equation.

One way to find the other solution is to approach this situation from underdamped case as a limit, i.e.
\[
\omega \to 0 \text{ from } \omega > 0
\]

\[
\begin{gathered}
	\begin{aligned}
		& \lim_{\omega \to 0} e^{ \frac{-\Gamma}{2}(t) } c\cos{(\omega t)} = e^{ \frac{-\Gamma}{2} t} \\
		& \lim_{\omega\to 0} e^{\frac{-\Gamma}{2} t} d \sin{\omega t} = 0
	\end{aligned}
\end{gathered}
\]

Divide the second term by $\omega$, because $\omega$ doesn't depend on $\omega$, $\frac{1}{\omega} e^{\frac{-\Gamma}{2} t} d\sin{\omega t}$ is still a solution.

\[
\begin{gathered}
	\lim_{\omega \to 0} e^{ \frac{-\Gamma}{2} t} \frac{ \sin{ \omega t} }{\omega} = te^{\frac{-\Gamma}{2} t} (\text{by L'Hopital's rule})
\end{gathered}
\]
So
\[
x(t) = ce^{ - \Gamma t/2} + dte^{-\Gamma t/2}
\]

\subsubsection{Forced oscillations} See Sec. 2.2 "Forced Oscillations" of Georgi (1992) \cite{agat:Geor1992}.
Equation of motion is
\begin{equation}\label{agat:Eq:ForcedOscillationEquationOfMotion:2}
	\ddot{x}(t) + \Gamma \dot{x} + \omega_0^2 x = F(t) /m 
\end{equation}
where
\begin{equation}
	F(t) = F_0 \cos{(\omega_dt)}
\end{equation}
See (2.1.5), (2.1.4) of Georgi (1992) \cite{agat:Geor1992}, respectively. \\

$\frac{\omega_d}{2\pi} = $ driving frequency, driving frequency can be tuned completely independently of the other parameters of the system.

If we let 
\[
\ddot{z}(t) + \Gamma \dot{z}(t) + \omega_0^2 z(t) = \mathcal{F}(t) /m
\]
where $\mathcal{F}(t) = F_0 e^{-i \omega_0 t}$.

For in particular numerical computation applications, we can transform Eq. \ref{agat:Eq:ForcedOscillationEquationOfMotion:2} into a system of first order equations: \\
Let 
\[
\begin{aligned}
& x_1(t) = x(t) \\
& x_2(t) = \dot{x}(t)
\end{aligned}
\]
so that
\[
\begin{aligned}
& \dot{x}_1 = x_2 \\
& \dot{x}_2 = -\Gamma x_2 - \omega_0^2 x_1 + f(t)
\end{aligned}
\]


Then if we look for steady state solution of the form
\[
z(t) = A e^{-i \omega_d t}
\]

One can think like the steady state solution is such that the system is driven with driving frequency $\omega_d$.

\begin{equation}
	(-\omega_d^2 + (-i \omega_d)\Gamma + \omega_0^2)A = \frac{F_0}{m} 
\end{equation}
See (2.20) of Georgi (1992) \cite{agat:Geor1992}

or
\begin{equation}
	A = \frac{ F_0 / m }{ \omega_0^2 - i\omega_d \Gamma - \omega_d^2}
\end{equation}
See (2.21) of Georgi (1992) \cite{agat:Geor1992}

The amplitude $A$ of the displacement is proportional to amplitude of driving force; this is as expected from linearity.

\[
\begin{gathered}
	\begin{aligned}
		& A = \frac{ \frac{F_0}{m} (\omega_0^2 + i \omega_d \Gamma - \omega_d^2 )}{ (\omega_0^2 - \omega_d^2)^2 + (\omega_d \Gamma)^2 } \text{ with } & \\
		& \Re{A} = \frac{ (\omega_0^2 - \omega_d^2) F_0/m }{(\omega_0^2 - \omega_d^2)^2 + (\omega_d \Gamma)^2} & \Im{A} \equiv B = \frac{ \Gamma \omega_d F_0 /m }{ (\omega_0^2 - \omega_d^2)^2 + (\Gamma \omega_d )^2 }
			\end{aligned} \\
		\mathcal{A} e^{- i\omega_d t} = \mathcal{A} ( c(\omega_dt) - i s(\omega_dt)) \xrightarrow{\Re} \Re{A } c(\omega_d t) + \Im{A} s(\omega_d t)
\end{gathered}
\]
where I use $\mathcal{A} \equiv A$ notation interchangeably, relying on context on whether it's real or complex, and $\cos{\omega_d t} \equiv c(\omega_d t)$, $\sin{\omega_d t} \equiv s(\omega_d t)$ as shorthand notation for sin and cos. \\

For 
\[
x(t) = \Re{z(t)} = (\Re{A}) c(\omega_d t) + Bs(\omega_d t)
\]

Term proportional to $B$ is $90^{\circ}$ out of phase. \\
Real part of $\mathcal{A}$, $A = \Re{A}$ is the \emph{elastic amplitude.} \\
Imaginary part of $\mathcal{A}$, $B = \Im{A}$ is called the \emph{absorptive amplitude}. \\

In the case where 
\[
	F(t) = F_0 \sin{(\omega_dt)}
\]
we'll see that the fundamental physical characteristics don't change, but the initial conditions for the solution do change, which is important for actual modeling (but gets looked over a bit in physics that concern about high-level concepts).

If $F(t) = F_0 \sin{(\omega_d t)}$, while the following doesn't change:
\[
\begin{gathered}
\begin{aligned}
& z(t) = \mathcal{A} e^{-i\omega_d t} \\
& \dot{z}(t) = (-i \omega_d) z \\
& \ddot{z}(t) = (-\omega_d^2)z
\end{aligned}
\end{gathered}
\]
but now consider 
\[
\begin{gathered}
\mathcal{F}(t) = + i F_0 e^{-i\omega_d t} = +i F_0 (c(\omega_d t)  -i s(\omega_d t) ) \\
\Re{(\mathcal{F}(t) )} = F_0 s(\omega_dt)
\end{gathered}
\]
so then
\[
\begin{gathered}
	(-\omega_d^2 + (-i \omega_d)\Gamma + \omega_0^2)\mathcal{A} = \frac{iF_0}{m}  \text{ or } 	\\
	\mathcal{A} = \frac{ iF_0 / m }{ \omega_0^2 - \omega_d^2- i\omega_d \Gamma} = \frac{ \frac{iF_0}{m} (\omega_0^2 - \omega_d^2 + i \omega_d \Gamma  )}{ (\omega_0^2 - \omega_d^2)^2 + (\omega_d \Gamma)^2 }
\end{gathered}
\]

So then in this case, 
\begin{equation}\label{agat:Eq:AmplitudeValuesForcedOscillationsForceProportionalToSine:2}
\begin{aligned}
& \Re{A} = \frac{ -\Gamma \omega_d F_0 /m }{ (\omega_0^2 - \omega_d^2)^2 + (\Gamma \omega_d )^2 } \equiv A \\
& \Im{A} = \frac{ (\omega_0^2 - \omega_d^2) F_0/m }{(\omega_0^2 - \omega_d^2)^2 + (\omega_d \Gamma)^2} \equiv B 
\end{aligned}
\end{equation}

and so the general form of the steady-state solution doesn't change, but the values for the amplitudes get switched:

\[
\begin{gathered}
x(t) = \Re{ z(t)} = Ac(\omega_d t) + Bs(\omega_dt) \\
\dot{x}(t) = \omega_d \left[ -A s(\omega_dt) + Bc(\omega_dt) \right]
\end{gathered}
\]

\subsubsection{Resonances} See Sec. 2.3 "Resonances" of Georgi (1992) \cite{agat:Geor1992}.

\[
\omega_d = \omega_0
\]
so $\Re{A} = 0$, $B = \frac{ F_0 / m}{ \Gamma \omega_d}$.

\subsubsection{Work, Work done by forced oscillation on the system}

Work is a nonlinear function of force, so we must use real force and real displacement.

\[
P(t) = F(t) \frac{\partial }{\partial t} X(t) = F_0 c(\omega_d t) (\omega_d) (-\Re{A} s(\omega_d t) + B c(\omega_d t))
\]

The first term is proportional to $s(2\omega_d t)$ and averages to 0 over a half period.

\[
\int_{t_0}^{t_0 + \pi / \omega_d} s(2\omega_d t) dt = \left. \frac{ -c(2\omega_d t) }{2 \omega_d } \right|_{t_0}^{t_0 + \pi /\omega_d} = 0
\]

This is why $\Re{A}$ is called elastic amplitude. If $A$ dominates, then energy fed into the system at 1 time is returned at a later time, as in elastic collision in mechanics.
\[
\begin{gathered}
	P_{\text{average}} \text{of 2nd. term} = P_{\text{average}} \\
	\text{ Consider } \\
	F_0 \omega_d B \int_{t_0}^{t_0 + \pi/ \omega_d} c^2 (\omega_d t) dt = F_0 \omega_d B \int_{t_0}^{t_0 + \pi/ \omega_d} \frac{ c(2\omega_d t) + 1}{2} dt = \\
	F_0 \omega_d B \left. \left( \frac{ s(2\omega_d t) }{ 4 \omega_d } + \frac{t}{2} \right) \right|_{t_0}^{t_0 + \frac{\pi}{\omega_d} } = F_0 \omega_d B \frac{\pi}{2\omega_d} = \frac{F_0 B\pi}{2} \\
	P_{\text{average}} = \frac{ \int_{t_0}^{t_0 + \pi / \omega_d} P(t) }{ \Delta t}  \text{ where } \Delta t = \pi / \omega_d \\
	\Longrightarrow P_{\text{average}} = \frac{ F_0 B\pi /2}{ \pi / \omega_d}  = \boxed{ \frac{F_0 \omega_d B}{2} }
\end{gathered}
\]

Recall that
\[
B = \frac{ \Gamma \omega_d F_0 /m }{ (\omega_0^2 - \omega_d^2)^2 + (\Gamma \omega_d)^2 }
\]

\[
\begin{gathered}
	P_{\text{average}} = \frac{F_0 \omega_d B}{2} = \frac{ (F_0 \omega_d)^2 \Gamma /2m }{ (\omega_0^2 - \omega_d^2)^2 + (\Gamma \omega_d)^2 } \xrightarrow{ \omega_d\to \omega_0} \frac{F_0^2}{ 2m \Gamma}
\end{gathered}
\]

The maximum average power is inversely proportional to $\Gamma$. \\

This analysis doesn't change when $F(t) = F_0 \sin{(\omega_d t)}$.

\[
\begin{gathered}
P(t) = F(t) \dot{x}(t) = F_0 s(\omega_d t) \omega_d [ -A s(\omega_d t) + Bc(\omega_d t) ] \\
\begin{aligned}
& \int_{t_0}^{t_0 + \frac{\pi}{\omega_d} } dt c(\omega_dt) s(\omega_d t) = \left. \frac{1}{2 \omega_d} s^2(\omega_d t) \right|_{t_0}^{t_0 + \frac{\pi}{\omega_d} } = 0 \\
& \int_{t_0}^{t_0 + \frac{\pi}{\omega_d} } dt s^2(\omega_dt) = \int_{t_0}^{t_0 + \frac{\pi}{\omega_d} } dt \frac{1 - c(2\omega_dt)}{2} = \frac{\pi}{2 \omega_d} - \left. \frac{s(2\omega_dt)}{ 4\omega_d} \right|_{t_0}^{t_0 + \frac{\pi}{\omega_d} } = \frac{\pi}{2\omega_d}
\end{aligned} \\
P_{\text{average}} = \frac{ \int_{t_0}^{t_0 + \frac{\pi}{\omega_d} } dt P(t) }{ \pi /\omega_d} = \frac{-F_0 A \pi}{2 / \pi/\omega_d} = \frac{-F_0 A \omega_d}{2}
\end{gathered}
\]
So $B$ is now the elastic amplitude and $A$ is the absorptive amplitude, but the same values are obtained, since $A$ and $B$ are effectively switched for $F(t) = F_0 \sin{\omega_d t}$.


To find the width (for fixed height) of the resonance curve (curve of $P_{\text{average}}$, depending on $\omega_d$), \\
find the value $\omega_d = \omega_{1/2}$ s.t. 
\[
P_{\text{average}} = \frac{P_{\text{average,max}}}{2} = \frac{F_0^2}{4m \Gamma}
\]
Then

\[
\begin{gathered}
\frac{P_{\text{average,max}}}{2} = \frac{F_0^2}{4m \Gamma} = \frac{ (F_0^2 \Gamma /2m ) \omega_{1/2}^2 }{ (\omega_{1/2}^2 - \omega_0^2)^2 + (\Gamma \omega_{1/2})^2 } \text{ or } \omega_{1/2}^4 -2\omega_{1/2}^2 \omega_0^2 + \omega_0^2 + \Gamma^2 \omega_{1/2}^2 = 2\Gamma^2 \omega_{1/2}^2 \text{ or } \\
\begin{aligned}
& \omega_{1/2}^4 - (2\omega_0^2 + \Gamma^2) \omega_{1/2}^2 + \omega_0^4 = 0 \\
& \omega_{1/2}^2 = \frac{ (2 \omega_0^2 + \Gamma^2) \pm \sqrt{ ( \Gamma^2 + 2\omega_0^2)^2 - 4\omega_0^4 } }{2} = \omega_0^2 + \frac{ \Gamma^2}{2} \pm \sqrt{ \frac{\Gamma^4 }{4}  + \Gamma^2 \omega_0^2 } = \omega_0^2 + \frac{ \Gamma^2}{2} \pm \Gamma \sqrt{ \frac{\Gamma^2 }{4}  +  \omega_0^2 }
\end{aligned}
\end{gathered}
\]
Now (EY: it seems like you would have needed to have known the answer beforehand to have tried this factorization)
\[
\begin{gathered}
\omega_{1/2}^2 = \left( \sqrt{ \omega_0^2 + \frac{\Gamma^2}{4} } \right)^2 \pm 2 \frac{\Gamma}{2} \sqrt{ \frac{\Gamma^2}{4} + \omega_0^2 } + \left( \frac{\Gamma}{2} \right)^2  = \left( \sqrt{ \omega_0^2 + \frac{ \Gamma^2}{4} } \pm \frac{\Gamma}{2} \right)^2
\end{gathered}
\]
Then
\begin{equation}
\boxed{ \omega_{1/2} = \sqrt{ \omega_0^2 + \frac{\Gamma^2}{4} } \pm \frac{\Gamma}{2} }
\end{equation}

The "free oscillation" of the system would be the case of underdamped oscillations; recall
\[
x(t) = A \exp{ \left( \frac{-\Gamma t}{2} \right) }\cos{ (\omega t - \theta) }
\]
Notice the exponential decay behavior of the amplitude or i.e. magnitude of $x(t)$ due to the $\exp{ \left( \frac{-\Gamma t}{2} \right)}$ factor. Let's just say the "lifetime" of free oscillation is the time for the amplitude to drop to half its value, i.e.
\[
\begin{gathered}
\frac{1}{2} = \exp{ \left( \frac{-\Gamma t_{1/2} }{2} \right) } \text{ or } -\ln{2} = \frac{-\Gamma t_{1/2} }{2} , \quad \text{ or } t_{1/2} = \frac{2\ln{2} }{\Gamma}
\end{gathered}
\]

For the time dependence of the free oscillation, i.e. $x(t) = A \exp{ \left( \frac{-\Gamma t}{2} \right) }\cos{ (\omega t - \theta) }$, then we say that the lifetime of the state in free oscillation is of order $1/\Gamma$, i.e.
\[
t_{1/2} = \frac{2\ln{2}}{ \Gamma}
\]

The width of the resonance peak in forced oscillation is 
\[
\Delta \omega_{1/2} = \Gamma
\]
so width of the resonance peak in forced oscillation is inversely proportional to the lifetime of the corresponding normal mode of free oscillation.

\subsection{Apply Fourier Transform to Forced Oscillations}

So for the equation of motion, Eq. \ref{agat:Eq:ForcedOscillationEquationOfMotion:2},
\[
\ddot{x} + \Gamma \dot{x} + \omega_0^2 x = F(t)
\]

Assume we drive the system at frequency $\omega_d$ and that system's response's frequency is also $\omega_d$.

Then

\[
\begin{gathered}
(i\omega_d)^2 \widehat{x} + (i\omega_d) \Gamma \widehat{x} + \omega_0^2 \widehat{x} = \widehat{F}(\omega_d) \text{ or } \frac{ \widehat{x}(\omega_d) }{ \widehat{F}(\omega_d) } = \frac{1}{ - (\omega_d^2 - i \omega_d \Gamma - \omega_0^2) }
\end{gathered}
\]

So in general, $\frac{ \widehat{x}(\omega_d) }{ \widehat{F}(\omega_d) }$ has a real part and imaginary part.

\subsection{Apply Laplace Transform to Forced Oscillations}

Recall Eq. \ref{agat:Eq:ForcedOscillationEquationOfMotion:2},
\[
\ddot{x} + \Gamma \dot{x} + k x = mF(t)
\]
so that if $\frac{k}{m} = \omega_0^2$, then
\[
\begin{gathered}
\xrightarrow{ \mathcal{L}} s^2 \widehat{x}(s) - s x(0) - x'(0) + \Gamma (s\widehat{x}(s) - x(0)) + \omega_0^2 \widehat{x}(s) = \widehat{F}(s) \text{ or } \\
(s^2 + \Gamma s + \omega_0^2) \widehat{x}(s) - (\Gamma + s) x(0) - x'(0) = \widehat{F}(s)
\end{gathered}
\]

If $x'(0) = 0$, and we start system at $x(0) =0$, so that we rely upon $F(t)$ to drive system, then we get a ratio of output or "response" to input, the driving, forced oscillation, per mass,
\[
\begin{gathered}
	\frac{ \widehat{x}(s) }{ \widehat{F}(s) } = \frac{1}{ s^2 + \Gamma s + \omega_0^2 } = \frac{1}{ \left( s + \frac{ \Gamma + \sqrt{ \Gamma^2 - 4 \omega_0^2 } }{ 2 } \right) \left( s + \frac{ \Gamma - \sqrt{ \Gamma^2 - 4 \omega_0^2 } }{ 2} \right) }
\end{gathered}
\]
Now look at where the denominator "blows up", the singularities. We had the case that if $\Gamma^2 - 4\omega^2 \geq 0 $ or $\frac{\Gamma}{2} \geq \omega_0$, we had overdamped or critically damped exponential decay. Then, the value of $s$ where a singularity occurs would be completely real. For the underdamped case with free oscillation, the value of $s$ where a singularity occurs would have an imaginary part.

If initial conditions can't be ignored, \\
then if $x'(0) =0$,
\[
\begin{gathered}
(s^2 + \Gamma s + \omega_0^2) \widehat{x}(s) - x'(0) = \widehat{F}(s) \text{ or } \frac{ \widehat{x}(s) }{ \widehat{F}(s) } = \frac{ 1 + x'(0) / \widehat{F}(s) }{ s^2 + \Gamma s + \omega_0^2 }
\end{gathered}
\]

But if $x'(0) =0$,
\[
\begin{gathered}
\frac{ \widehat{x}(s) }{ \widehat{F}(s) } = \frac{ 1 + (\Gamma +s )x(0) / \widehat{F}(s) }{ s^2 + \Gamma s + \omega_0^2 }
\end{gathered}
\]

\subsubsection{General solution to forced oscillation}

For subscript $s$ meaning steady-state driven solution, see Eq. \ref{agat:Eq:ForcedOscillationSteadyStateSolutionAmplitude}
\[
\begin{gathered}
x_s(t) = Ac(\omega_d t) + Bs(\omega_d t) \\
\begin{aligned}
& A = \frac{ (\omega_0^2 - \omega_d^2) F_0 / m }{ (\omega_0^2 - \omega_d^2)^2 + (\omega_d \Gamma)^2 } \\
& B = \frac{ \Gamma \omega_d F_0 / m }{ (\omega_0^2 - \omega_d^2)^2 + (\Gamma \omega_d)^2 }
\end{aligned}
\end{gathered}
\] 
so then
\[
\begin{aligned}
& x_s(0) = A \\
& \dot{x}_s(0) = \omega_dB
\end{aligned}
\]

Let $x_{\text{free}}$ denote the response, or solution, or the resulting motion for the "free oscillation", the underdamped, exponential decaying solution, for $\omega_0 > \frac{\Gamma}{2}$. See Eq. \ref{agat:Eq:UnderdampedOscillationsSolution:2},

\[
\begin{gathered}
\begin{aligned}
x_{\text{free}}(t) & = \exp{ \left( \frac{-\Gamma t}{2} \right) } (c\cos{(\omega t)} + d\sin{ (\omega t)} ), \quad \, \omega^2 = \omega_0^2 -\frac{\Gamma^2}{2} \\
x_{\text{free}}(0) &= c \\
\dot{x}_{\text{free}}(t) & = \frac{-\Gamma}{2} x_{\text{free}}(t) + e^{-\frac{\Gamma}{2} t} \omega( - cs(\omega t) + dc(\omega t) ) \\
\dot{x}_{\text{free}}(0) & = \frac{-\Gamma}{2} c + \omega d
\end{aligned}
\end{gathered}
\]

The most general solution $x(t)$ is of the form
\[
\begin{gathered}
x(t) = x_{\text{free}}(t) + x_s(t)
\end{gathered}
\]
with initial conditions
\[
\begin{aligned}
& x(0) = c + A \text{ so } c = x(0) - A \\
& \dot{x}(0) = \frac{-\Gamma}{2} c + \omega d + \omega_d B \text{ so } d = \frac{1}{\omega} \left[ \dot{x}(0) + \frac{\Gamma}{2} (x(0) - A) - \omega_d B \right]
\end{aligned}
\]

And so the full, general solution is 
\begin{equation}
\begin{gathered}
x(t) = e^{\frac{-\Gamma t}{2} } (c\cos{(\omega t)} + d \sin{ (\omega t)} ) + A \cos{ (\omega_d t)} + B \sin{(\omega_d t)} \\
\text{ where } \\
\begin{aligned}
& A = \frac{ (\omega_0^2 - \omega_d^2) F_0 / m }{ (\omega_0^2 - \omega_d^2)^2 + (\omega_d \Gamma)^2 } \\
& B = \frac{ \Gamma \omega_d F_0 / m }{ (\omega_0^2 - \omega_d^2)^2 + (\Gamma \omega_d)^2 }
& \omega^2 = \omega_0^2 - \frac{\Gamma^2}{4}, \\
& c = x(0) - A, \\
& d = \frac{1}{\omega} \left[ \dot{x}(0) + \frac{\Gamma}{2} (x(0) - A) - \omega_d B \right]
\end{aligned}
\end{gathered}
\end{equation}

For the case of $F(t) = F_0 \sin{(\omega_d t)}$, i.e. the force is proportional to $\sin$, then the exact same analysis leads to the same result as above, but for $A,B$ values switched to be as Eq. \ref{agat:Eq:AmplitudeValuesForcedOscillationsForceProportionalToSine:2}.

\[
\begin{aligned}
& \Re{A} = \frac{ -\Gamma \omega_d F_0 /m }{ (\omega_0^2 - \omega_d^2)^2 + (\Gamma \omega_d )^2 } \equiv A \\
& \Im{A} = \frac{ (\omega_0^2 - \omega_d^2) F_0/m }{(\omega_0^2 - \omega_d^2)^2 + (\omega_d \Gamma)^2} \equiv B 
\end{aligned}
\]

